\section{Distorsione e Reprojection Error}
\label{sec:distorsione}

\subsection{Modello di distorsione radiale e tangenziale}

Il modello pinhole suppone una trasformazione lineare, ma le lenti reali
introducono \emph{distorsione}: il punto effettivamente osservato
\((x_d,y_d)\) differisce dal punto teorico \((x_n,y_n)\) secondo
\begin{align}
  x_d &= x_n(1+k_1r^2+k_2r^4+k_3r^6)
        +2p_1x_ny_n+p_2(r^2+2x_n^2),\\
  y_d &= y_n(1+k_1r^2+k_2r^4+k_3r^6)
        +p_1(r^2+2y_n^2)+2p_2x_ny_n,
\end{align}
dove \(r^2=x_n^2+y_n^2\), \(k_1,k_2,k_3\) sono i coefficienti di
distorsione radiale e \(p_1,p_2\) quelli
tangenziali~\cite[Sez.~6.5]{HZ2004}.

\subsection{Reprojection error --- la regola aurea}
\label{sub:reproj}

Una volta stimati \(\mat{K}\), \(\mathbf{k}_r=(k_1,k_2,p_1,p_2,k_3)\)
e gli estrinseci \(\{[\mat{R}_j|\mathbf{t}_j]\}\), la qualità della
calibrazione si misura con il \emph{reprojection error} medio:
\[
  \epsilon_{\text{RMSE}}
  =\sqrt{\frac{1}{N}\sum_{j=1}^{n}\sum_{i=1}^{m_j}
          \norm{\homo{x}_{ji}-\hat{\homo{x}}_{ji}}^2},
  \label{eq:reproj}
\]
dove \(\hat{\homo{x}}_{ji}=\pi(\mat{K}[\mat{R}_j|\mathbf{t}_j]\homo{X}_i,
\mathbf{k}_r)\) è la riproiezione del punto 3D~\cite[Sez.~5.1]{HZ2004}.

Questa distanza nel piano immagine è la \emph{Gold Standard}: minimizzare
\(\epsilon_{\text{RMSE}}\) equivale a trovare i parametri che meglio
spiegano le corrispondenze osservate~\cite[Sez.~5.1]{HZ2004}.

\subsection{Motivazione per il GA}

Il raffinamento finale di Zhang minimizza \eqref{eq:reproj} tramite
Levenberg-Marquardt (LM), richiedendo il Jacobiano della funzione di
riproiezione. LM è localmente efficiente, ma dipende dalla qualità delle
immagini: immagini sfocate, pose troppo simili o corrispondenze errate
degradano l'ottimo locale raggiunto.

La \emph{selezione del sottoinsieme} di immagini è un problema
combinatorio su cui LM non agisce: è esattamente questo spazio che
l'Algoritmo Genetico esplora (Sezione~\ref{sec:ga}).